\bookchapter{Глава 47: Справочник и ресурсы}

\chapterintro{Эта глава — справочник, которым можно пользоваться и отдельно. Когда забывается какой-то принцип или нужно быстро принять решение, открывайте прямо здесь. Советуем держать его как «карманные карточки»: шпаргалка пяти принципов, дерево выбора инструментов, частые вопросы, ресурсы для обучения — ничего не пропало.}

\booksection{Карточки шпаргалки мыслительных моделей}

\begin{center}
\begin{tabularx}{\linewidth}{|X|X|X|}
\hline
\textbf{Принцип} & \textbf{Одной фразой} & \textbf{Ключ к действию} \tabularnewline
\hline
Think Before Coding & Сначала подумай — потом спрашивай & Ответьте на три вопроса: что делаем? как выглядит результат? какие ограничения? \tabularnewline
Simplicity First & Решай задачу минимальным кодом & За один раз — одна функция, отсекай лишнюю сложность \tabularnewline
Surgical Changes & Правь точно и не выходи за рамки & Явно указывайте область изменений и отклоняйте «заодно поправил» от ИИ \tabularnewline
Goal-Driven Execution & Говори «что сделать», а не «как делать» & Давайте цель, а не шаги, — пусть ИИ сам выбирает путь \tabularnewline
Iterate \& Verify & Цикл «сгенерировал — проверил — исправил» & Трёхуровневая проверка: функциональность, границы, архитектура \tabularnewline
\hline
\end{tabularx}
\end{center}

\begin{center}\small\color{gray}Шпаргалка пяти принципов мышления\end{center}

\booksection{Дерево выбора инструментов}

Когда нужно выбрать инструмент для программирования с ИИ, идите по следующему дереву решений. Важно: это не рейтинг «что лучше», а рамка для ответа «что подходит именно вам».

Что вы хотите сделать?

С нуля, хочется быстро попробовать выбирайте облачный Builder (v0/Bolt/Replit)

Повседневная разработка, нужно автодополнение выбирайте IDE со встроенным ИИ (Cursor/Trae)

Крупные проекты, нужны массовые операции выбирайте терминальный агент (Claude Code)

Проверить идею быстрым прототипом облачный Builder или Trae SOLO

\begin{center}\small\color{gray}Рисунок 6: дерево выбора инструмента для программирования с ИИ\end{center}

\booksection{Частые вопросы (FAQ)}

В 1: Я совсем не умею программировать. Смогу ли я заниматься программированием с ИИ?

Сможете. К 2026 году инструменты программирования с ИИ уже настолько сильны, что можно описывать требования обычным языком и получать от ИИ код. Но два совета: во-первых, начните с облачных Builder-инструментов (вроде v0, Bolt) — они не требуют настройки среды разработки; во-вторых, по ходу дела осваивайте базовые понятия — знать синтаксис в совершенстве не нужно, но надо уметь «улавливать общий смысл» кода, который пишет ИИ, иначе грамотно принять работу не получится.

В 2: Какой язык программирования мне стоит учить?

Зависит от того, что вы хотите делать. Фронтенд сайтов — JavaScript/TypeScript; обработка данных или ИИ — Python; корпоративный бэкенд — Java или Go. Но важнее другое: в эпоху программирования с ИИ значимость «языка» снижается. Синтаксическими деталями займётся ИИ, а вам нужно понимать «структуру программы» — что такое функция, что такое цикл, что такое структуры данных. Эти понятия межъязыковые: поняв их, легко перейдёте на любой язык.

В 3: Есть ли проблемы с авторскими правами на код, написанный ИИ?

По состоянию на август 2026 года это всё ещё спорный юридический вопрос. В целом код, сгенерированный ИИ, в большинстве юрисдикций не считается имеющим определённого правообладателя. Но если вы используете его в коммерческом проекте, советуем: 1) изучить в условиях службы вашего ИИ-инструмента пункты об интеллектуальной собственности; 2) вручную проверять и дорабатывать ключевой код, внося собственный творческий вклад; 3) не позволять ИИ напрямую копировать открытый код, который может охраняться авторским правом.

В 4: Заменит ли ИИ программистов?

В ближайшей перспективе — нет. ИИ вытесняет не профессию «программист», а занятие «кодирование построчно». Точно так же калькулятор не заменил математиков, а освободил их от утомительных ручных вычислений. Инструменты программирования с ИИ переводят программиста от «писать код» к «думать над задачами» — это на деле поднимает потолок ценности программиста. Конкурентоспособность будущего определяется не тем, «сколько строк кода можешь написать», а тем, «насколько сложные задачи можешь решать».

В 5: Как понять, правильно ли то, что выдал ИИ?

Метод трёхуровневой проверки. Первый уровень — функциональная проверка: запускается ли? верны ли вход и выход? Это самая база. Второй уровень — проверка границ: что будет при пустом вводе, сверхдлинном вводе, некорректном вводе? Так вскрываются скрытые проблемы. Третий уровень — архитектурная проверка: разумна ли структура кода? нет ли зашитых в код значений? удобно ли расширять? Здесь нужен уже технический кругозор. Если третий уровень пока не по силам — делайте хотя бы первые два: «запускается и границы в порядке» уже лучше, чем у большинства новичков.

\booksection{Рекомендуемые ресурсы для обучения}

Ниже — рекомендуемые ресурсы по состоянию на август 2026 года. Учтите: ссылки в интернете могут меняться; если какая-то не работает, по названию обычно легко найти актуальную версию.

\begin{center}
\begin{tabularx}{\linewidth}{|X|X|X|}
\hline
\textbf{Тип} & \textbf{Ресурс} & \textbf{Почему рекомендуем} \tabularnewline
\hline
Первоисточник & CLAUDE.md от Andrej Karpathy & Источник пяти принципов: 70 строк, но каждая строка — золото \tabularnewline
Первоисточник & Официальная документация Anthropic (Claude Code) & Лучшие практики для терминального агента \tabularnewline
Уроки & Сборник Skills от сообщества Vibe Coding & Набор практических приёмов от участников сообщества \tabularnewline
Уроки & Официальная документация и YouTube-канал Cursor & Авторитетное руководство по ИИ, встроенному в IDE \tabularnewline
Сообщество & awesome-claude-code на GitHub & Лучшие практики и набор плагинов для Claude Code \tabularnewline
Сообщество & Сообщество r/cursor на Reddit & Обмен опытом и обсуждения пользователей Cursor со всего мира \tabularnewline
База & MDN Web Docs (developer.mozilla.org) & Самый авторитетный справочник по фронтенд-разработке \tabularnewline
База & Официальное руководство по Python (docs.python.org) & Официальная стартовая точка изучения Python \tabularnewline
\hline
\end{tabularx}
\end{center}

\begin{center}\small\color{gray}Список рекомендуемых ресурсов для обучения\end{center}

\booksection{Заключение: инструменты в руках, методология в голове}

Вернёмся к главному убеждению книги: инструменты обновляются каждые три месяца, методология эволюционирует раз в три года. Сегодня, в августе 2026 года, инструменты программирования с ИИ уже поразительно сильны: одной фразой вы создаёте сайт, одним описанием — fullstack-приложение, одной инструкцией — сложный рефакторинг. Но чем сильнее инструменты, тем заметнее важность методологии — сила инструментов увеличивает вес каждого вашего решения.

Если ваши требования ясны, ИИ поможет быстро их воплотить; если размыты — ИИ так же быстро сделает неправильно. Если вы следуете принципу «простота прежде всего», ИИ поможет быстро проверить идею; если гнаться за всем сразу — ИИ так же быстро нагребёт технического долга. Если вы правите точно, ИИ поможет хирургически отшлифовать код; если инструкция мутная — ИИ размашисто насыплет багов.

Именно поэтому книга посвятила большую часть методологии, а не оперированию инструментами. Не потому, что инструменты не важны, а потому, что методология — «множитель» инструментов: чем лучше методология, тем выше ценность инструментов; при нулевой методологии самый мощный инструмент — просто ноль, умноженный на любое число.

Вы больше не программируете в одиночку. Вы — командир, а ИИ — ваша армия. Качество вашего мышления определяет боевую мощь армии. С сегодняшнего дня вооружайте мышление методологией и усиливайте возможности инструментами — в этом весь секрет программирования с ИИ.

\booksection{Список литературы}

\begin{enumerate}
  \item Andrej Karpathy. (2025). Vibe Coding. Social media post introducing the concept of vibe-driven programming with AI.
  \item Andrej Karpathy. (2025). CLAUDE.md - Project guidelines for AI-assisted development. GitHub repository, 70+K stars.
  \item Anthropic. (2026). Claude Code: Terminal-native AI programming. Official documentation and best practices.
  \item Cursor Team. (2026). Cursor 3: The next generation of AI-first IDE. Release notes and documentation.
  \item OpenAI. (2026). Codex CLI: Bringing AI to the terminal. Technical documentation.
  \item Model Context Protocol (MCP) Specification. (2025-2026). Open protocol for AI-tool integration.
  \item Vibe Coding Community. (2026). Community-curated skills and best practices for AI-assisted programming.
  \item Spec-Driven Development (SDD) methodology. (2026). Emerging practices in specification-first AI programming.
\end{enumerate}
